\documentclass[11pt]{letter}
\usepackage[utf8]{inputenc}
\usepackage[T1]{fontenc}
\usepackage[margin=1in]{geometry}
\usepackage{hyperref}
\usepackage{siunitx}
\urlstyle{same}

\signature{Jonathan Mace\\Corresponding author}
\address{Jonathan Mace\\Microsoft Research\\Redmond, WA\\jonathanmace@microsoft.com}
\date{\today}

\begin{document}

\begin{letter}{%
Editor-in-Chief\\
\textit{Journal of the Acoustical Society of America}}

\opening{Dear Editor-in-Chief,}
\small

Please consider our manuscript, ``Can You Feel the Bass? Quantitative Partition of Airborne and Structure-Borne Pathways for Sub-Bass Perception at Concert Sound Pressure Levels,'' for publication in \textit{Journal of the Acoustical Society of America}. The paper addresses a familiar acoustical claim --- that intense low-frequency sound can be felt viscerally --- by quantitatively separating the airborne and structure-borne pathways that contribute to this sensation. Extending our fluid-filled viscoelastic shell model of the abdomen into the \SIrange{20}{80}{\hertz} sub-bass band, we use an energy-consistent reciprocity analysis for airborne coupling and a damped-plate power-balance model for floor and seating transmission. The result is a clear pathway partition: airborne sound is a genuine near-miss rather than utterly negligible, but structure-borne vibration remains the dominant and perceptible route at concert levels.

We believe the manuscript fits \textit{JASA} because it combines acoustics, perception, and biomedical relevance in a single physically interpretable framework. The paper replaces pressure-based airborne estimates with an energy-consistent treatment, compares predicted displacements against whole-body vibration perception thresholds, and links concert-acoustics conditions to organ-scale biomechanical response. It should therefore interest readers in low-frequency acoustics, concert and architectural sound, vibrotactile perception, and biomedical acoustics. The manuscript also forms part of our broader Browntone research programme on low-frequency acoustic and structure-borne stimulation of the human body.

The manuscript is original, has not been published previously, is not under consideration elsewhere, and has been approved by both authors. Brian R. Mace's submission email is \texttt{b.mace@auckland.ac.nz}. We hope the paper will be of interest to \textit{JASA} as a quantitative account of why listeners can indeed feel the bass, but mostly through the floor rather than through the air alone.

\closing{Yours sincerely,}

\end{letter}
\end{document}
